\documentclass[10pt,a4paper]{amsart}
\usepackage[margin=19mm]{geometry}
\usepackage{amsmath,amssymb,mathtools}
\usepackage[scr=rsfs,cal=euler]{mathalfa}
\usepackage{xfrac}
\usepackage[unicode,hidelinks,bookmarks=false]{hyperref}
\newcommand{\ie}{i.e.\ }
\newcommand{\cf}{cf.\ }
\newcommand{\resp}{respectively\ }
\newcommand{\Subsection}{\subsection}
\providecommand{\nobreakdash}{\nobreak\hbox{-}}
\setlength{\emergencystretch}{2em}
\multlinegap0pt
\usepackage[shortlabels]{enumitem}
\usepackage{amssymb}
\usepackage{mathtools}
\usepackage{xcolor}
\usepackage{tikz}
\newtheorem{lemma}{Lemma}
\newtheorem{proposition}{Proposition}
\newcommand{\Ebb}{\mathbb{E}}
\newcommand{\cL}{\mathcal{L}}
\newcommand{\dd}{\mathrm{d}}
\title{Macroscopic loops in the random loop model on sparse random graphs}
\author{Andreas Klippel}
\date{Source: arXiv:2604.20514v1 (CC BY 4.0)}
\begin{document}
\maketitle

\noindent\emph{Source and licence.} Macroscopic loops in the random loop model on sparse random graphs. Version arXiv:2604.20514v1 (CC BY 4.0), distributed by arXiv under the Creative Commons Attribution 4.0 International licence, http://creativecommons.org/licenses/by/4.0/, as recorded in the arXiv licence metadata for this exact version and shown on the abstract page https://arxiv.org/abs/2604.20514v1. Full licence notice: \texttt{licenses/arxiv-2604.20514-CC-BY-4.0.txt}.\par\smallskip

\noindent\textit{Editorial scope.} Counted unit: the article's proposition prop:exact-drift (source0.tex lines 994--1012). The article's complete original proof of it is delivered with it (source0.tex lines 1014--1192, one \texttt{proof} environment of 179 lines). No mathematics is written, repaired or re-derived: the statement and the proof are the article's own lines, in the article's own notation, and no retained line is substituted or re-flowed.  Retained uncounted as the source-local context the proof needs: lines 886--902; lines 933--945.  Nothing was omitted from any retained span, so the package carries no omission record.  No retained line contains a citation, so the wrapper carries no bibliography and no bibliography entry.  No retained line contains a figure environment, an image-inclusion command or drawing code, so no image is delivered and the package declares no asset dependency.  Editorial formatting only, in addition: an AMS \texttt{amsart} preamble that restates the article's own declarations and macros used by the retained lines, namely the environments \texttt{lemma}, \texttt{proposition} on separate counters, together with the standard packages those lines need.  The retained environments print in this document as Lemma 1, Proposition 1 in order of appearance, whereas the article prints the same items with the numbering of its own full document.  The complete native source of the article as distributed in the arXiv e-print for this exact version is bundled unchanged as \texttt{sources/198941/source0.tex}.
\begin{lemma}[Local split--merge--rewire rule]
\label{lem:local-surgery}
Let \(\omega\) be a configuration, let \(\cL(\omega)\) be the associated loop
configuration, let \(e=\{x,y\}\in E\), and let \(s\in S^1\setminus T(\omega)\) be a
regular time. Consider the effect of inserting at \((e,s)\) either a cross or a bar.

\begin{enumerate}[label=(\roman*)]
    \item If \(I_{\neq}(\omega;e,s)=1\), then inserting either a cross or a bar
    merges the two loops containing \((x,s)\) and \((y,s)\) into one.

    \item If \(I_+(\omega;e,s)=1\), then inserting a cross splits the loop into two,
    whereas inserting a bar rewires the loop without changing the number of loops.

    \item If \(I_-(\omega;e,s)=1\), then inserting a bar splits the loop into two,
    whereas inserting a cross rewires the loop without changing the number of loops.
\end{enumerate}
\end{lemma}

\begin{lemma}[Differentiation formula for a finite-intensity Poisson process]
\label{lem:poisson-diff}
Let \((\mathsf X,\nu)\) be a finite measure space, and for \(t\ge 0\) let \(\Pi_t\) be
a Poisson point process on \(\mathsf X\) with intensity measure \(t\nu\). Then, for
every finite counting measures, the map
\(t\mapsto \Ebb[F(\Pi_t)]\) is differentiable on \((0,\infty)\), and its right
derivative at \(t=0\) is given by
\[
\frac{\dd}{\dd t}\Ebb[F(\Pi_t)]
=
\int_{\mathsf X}\Ebb\!\left[F(\Pi_t+\delta_z)-F(\Pi_t)\right]\nu(\dd z).
\]
\end{lemma}

\begin{proposition}[Differential identity]
\label{prop:exact-drift}
Let \(G=(V,E)\) be a finite graph, let \(u\in[0,1]\), and let \(\theta>0\). Then
\[
D_{\theta,u}(t)
=
\Ebb_{\theta,t,u}^G\!\left[
(1+\theta u)J_+ + (1+\theta(1-u))J_- - |E|
\right],
\]
where
\[
D_{\theta,u}(t)=
\begin{cases}
\dfrac{\dd}{\dd t}\,\Ebb_{1,t,u}^G[\lambda(\omega)], & \theta=1,\\[1.2ex]
\dfrac{\theta}{\theta-1}\dfrac{\dd}{\dd t}\log Z_G(\theta,t,u), & \theta\neq 1.
\end{cases}
\]
\end{proposition}

\begin{proof}
We encode the model as a marked Poisson point process on
\[
\mathsf X:=E\times S^1\times\{\times,\mid\}.
\]
Let \(\nu_u\) be the finite measure on \(\mathsf X\) defined by
\[
\nu_u(\dd e,\dd s,\dd m)
:=
\sum_{e\in E}\dd s\,
\bigl(u\,\delta_{\times}+(1-u)\,\delta_{\mid}\bigr)(\dd m).
\]
Thus \(\nu_u(\mathsf X)=|E|\), and the underlying random loop configuration is a
Poisson point process \(\omega\) on \(\mathsf X\) with intensity measure \(t\nu_u\).
Its law is denoted by \(\rho_{t,u}^G\).

For \(z=(e,s,m)\in\mathsf X\), we write \(\omega+\delta_z\) for the configuration
obtained by inserting one additional marked point \(z\). Since
\(\lambda(\omega)\in\{1,\dots,|V|\}\), the functionals
\[
F_\theta(\omega):=\theta^{\lambda(\omega)}
\qquad\text{and}\qquad
F_1(\omega):=\lambda(\omega)
\]
are bounded.
We first treat the case \(\theta\neq 1\). By Lemma~\ref{lem:poisson-diff},
\[
\frac{\dd}{\dd t}Z_G(\theta,t,u)
=
\int_{\mathsf X}
\Ebb_{\rho_{t,u}^G}\!\left[
\theta^{\lambda(\omega+\delta_z)}-\theta^{\lambda(\omega)}
\right]\nu_u(\dd z).
\]
Write \(z=(e,s,m)\). Since the exceptional set \(T(\omega)\subset S^1\) is finite,
the identities from Lemma~\ref{lem:local-surgery} hold for \(\nu_u\)-almost every
\((e,s,m)\).
If \(m=\times\), then Lemma~\ref{lem:local-surgery} implies
\[
\lambda(\omega+\delta_{(e,s,\times)})-\lambda(\omega)
=
I_+(\omega;e,s)-I_{\neq}(\omega;e,s).
\]
Likewise, if \(m=\mid\), then
\[
\lambda(\omega+\delta_{(e,s,\mid)})-\lambda(\omega)
=
I_-(\omega;e,s)-I_{\neq}(\omega;e,s).
\]
Hence, for \(\nu_u\)-almost every \((e,s,m)\),
\[
\theta^{\lambda(\omega+\delta_{(e,s,\times)})}-\theta^{\lambda(\omega)}
=
\theta^{\lambda(\omega)}
\Bigl((\theta-1)I_+ + (\theta^{-1}-1)I_{\neq}\Bigr),
\]
and
\[
\theta^{\lambda(\omega+\delta_{(e,s,\mid)})}-\theta^{\lambda(\omega)}
=
\theta^{\lambda(\omega)}
\Bigl((\theta-1)I_- + (\theta^{-1}-1)I_{\neq}\Bigr).
\]
Integrating over the mark distribution
\(u\delta_\times+(1-u)\delta_{\mid}\), we get
\[
\int_{\{\times,\mid\}}
\bigl(\theta^{\lambda(\omega+\delta_{(e,s,m)})}-\theta^{\lambda(\omega)}\bigr)
\bigl(u\delta_\times+(1-u)\delta_{\mid}\bigr)(\dd m)
\]
\[
=
\theta^{\lambda(\omega)}
\Bigl(
(\theta-1)u\,I_+ + (\theta-1)(1-u)\,I_- + (\theta^{-1}-1)I_{\neq}
\Bigr).
\]
Now note that for almost every \((e,s)\),
\[
I_+(\omega;e,s)+I_-(\omega;e,s)+I_{\neq}(\omega;e,s)=1,
\]
because exactly one of the three alternatives above holds. Therefore
\[
\sum_{e\in E}\int_{S^1} I_{\neq}(\omega;e,s)\,\dd s
=
|E|-J_+(\omega)-J_-(\omega).
\]
Using also
\[
\sum_{e\in E}\int_{S^1} I_+(\omega;e,s)\,\dd s = J_+(\omega),
\qquad
\sum_{e\in E}\int_{S^1} I_-(\omega;e,s)\,\dd s = J_-(\omega),
\]
we obtain
\[
\frac{\dd}{\dd t}Z_G(\theta,t,u)
=
\Ebb_{\rho_{t,u}^G}\!\left[
\theta^{\lambda(\omega)}
\Bigl(
(\theta-1)u\,J_+
+
(\theta-1)(1-u)\,J_-
+
(\theta^{-1}-1)(|E|-J_+-J_-)
\Bigr)
\right].
\]
Since \(\theta^{-1}-1=-(\theta-1)/\theta\), the bracket equals
\[
\frac{\theta-1}{\theta}
\Bigl(
(1+\theta u)J_+ + (1+\theta(1-u))J_- - |E|
\Bigr).
\]
Therefore
\[
\frac{\dd}{\dd t}Z_G(\theta,t,u)
=
\frac{\theta-1}{\theta}\,
\Ebb_{\rho_{t,u}^G}\!\left[
\theta^{\lambda(\omega)}
\Bigl(
(1+\theta u)J_+ + (1+\theta(1-u))J_- - |E|
\Bigr)
\right].
\]
Dividing by \(Z_G(\theta,t,u)\) yields
\[
\frac{\theta}{\theta-1}\frac{\dd}{\dd t}\log Z_G(\theta,t,u)
=
\Ebb_{\theta,t,u}^G\!\left[
(1+\theta u)J_+ + (1+\theta(1-u))J_- - |E|
\right].
\]
For \(\theta=1\), apply Lemma~\ref{lem:poisson-diff} to the bounded functional
\(F(\omega)=\lambda(\omega)\). Then
\[
\frac{\dd}{\dd t}\Ebb_{1,t,u}^G[\lambda(\omega)]
=
\int_{\mathsf X}
\Ebb_{1,t,u}^G\!\left[
\lambda(\omega+\delta_z)-\lambda(\omega)
\right]\nu_u(\dd z).
\]
By Lemma~\ref{lem:local-surgery},
\[
\lambda(\omega+\delta_{(e,s,\times)})-\lambda(\omega)=I_+-I_{\neq},
\]
\[
\lambda(\omega+\delta_{(e,s,\mid)})-\lambda(\omega)=I_--I_{\neq}
\]
for \(\nu_u\)-almost every \((e,s)\). Averaging over the mark distribution gives
\[
u(I_+-I_{\neq})+(1-u)(I_--I_{\neq})
=
uI_+ + (1-u)I_- - I_{\neq}.
\]
Using \(I_{\neq}=1-I_+-I_-\), this becomes
\[
(1+u)I_+ + (2-u)I_- - 1.
\]
Integrating over \(e\) and \(s\) yields
\[
\frac{\dd}{\dd t}\Ebb_{1,t,u}^G[\lambda(\omega)]
=
\Ebb_{1,t,u}^G\!\left[
(1+u)J_+ + (2-u)J_- - |E|
\right].
\]
Since
\[
1+u=1+\theta u,
\qquad
2-u=1+\theta(1-u)
\qquad\text{when }\theta=1,
\]
this is exactly the claimed formula.
\end{proof}

\end{document}
